% Section 6: Discussion
% Note: \section{} header is in main.tex

\subsection{Practical Implications}

Existing diversity-preserving methods (DIVE \cite{placeholder_dive}, DARLING \cite{placeholder_darling}, GEM \cite{placeholder_gem}) target semantic diversity as their primary objective, assuming that reward optimization narrows the semantic distribution. Our results challenge this assumption: across eight configurations, Vendi Score remains within $\pm$1\% of its initial value. The collapse that does occur is structural. Opening entropy drops by 0.75 bits under both curated and random conditions at $K{=}4$, intensifying to 1.13 bits at $K{=}8$. The near-duplicate start rate rises from 37.9\% to 75.6\% regardless of curation strategy. A practitioner who monitors only Vendi Score or Self-BLEU will conclude diversity is intact, while end users encounter structurally uniform responses.

Users perceive repetition through surface patterns, not embedding geometry. We recommend that diversity audits include at least one format-level metric alongside embedding-based measures. Opening entropy and near-duplicate start rate require no GPU and capture homogenization that current pipelines miss.

A separate practical constraint concerns the reward signal itself. The three proxy dimensions used in evaluation (quality, diversity, complexity) correlate at $r = 0.63$--$0.77$, limiting the effectiveness of single-objective curation. This entanglement explains both the near-zero quality improvement across generations and, paradoxically, why curation does not collapse semantic diversity: a reward signal correlated with multiple text properties selects for a broader band of responses than one isolating a single dimension. If the same reward model that curates training data also serves as the quality metric, improvements are circular; independent evaluation through held-out judges or orthogonal metrics remains essential.

\subsection{Methodological Considerations}

The theoretical conditions for model collapse \cite{placeholder_shumailov} assume unfiltered self-consumption where all generated data is recycled. Reward-curated pipelines violate two preconditions: the training distribution is filtered by an external reward signal (not the model's full output), and the prompt set provides an external scaffold forcing broad semantic coverage. Reward saturation further limits compounding, with a score gap of only 0.04 between the top-1 candidate and pool average. These conditions are sufficient for semantic stability in our setting, but they are not guaranteed in all deployments. A narrow prompt domain would weaken the external scaffold, a highly discriminative reward model would sharpen selection pressure beyond the saturation regime we observe, and deterministic decoding would eliminate the variance injection that counteracts convergence. Practitioners operating in any of these regimes should verify semantic stability empirically rather than relying on our results.

The direction reversal between LoRA and Full-FT has broader methodological implications. Under LoRA, curation suppresses Vendi Score by 3.40 pp relative to random; under Full-FT at the same learning rate, the gap shrinks to 0.46 pp; at higher learning rate, curation promotes diversity by 0.63 pp. This progression is consistent with LoRA's rank-16 bottleneck imposing an upper bound on the distributional structure the model can absorb. Conclusions drawn from LoRA-based evaluations of curation may conflate capacity constraints with method properties. When the research question concerns distributional effects of data selection, evaluation under full fine-tuning is necessary to separate capacity effects from genuine curation dynamics.
